\section{Giới thiệu}

Trong xử lý ảnh kỹ thuật số, nhiễu (noise) là một vấn đề không thể tránh khỏi, phát sinh từ nhiều nguồn khác nhau như cảm biến quang học, quá trình truyền tín hiệu, hay điều kiện ánh sáng kém. Bài toán image denoising --- khôi phục ảnh gốc từ ảnh bị nhiễu --- đã và đang là một trong những bài toán cơ bản nhất của lĩnh vực image processing, với ứng dụng rộng rãi từ y tế, viễn thám, đến nhiếp ảnh số.

Một trong những hướng tiếp cận thống trị trong nhiều thập kỷ là transform-domain thresholding: phân tích ảnh sang một miền biến đổi, loại bỏ các hệ số nhỏ (được xem là noise), rồi tái tạo lại ảnh. Phương pháp này hoạt động hiệu quả dựa trên giả định rằng tín hiệu ảnh có biểu diễn sparse trong miền transform được chọn --- tức là hầu hết các hệ số gần bằng không, trong khi chỉ một số ít hệ số lớn mang thông tin tín hiệu. Wavelet transform, đặc biệt là Discrete Wavelet Transform (DWT) kết hợp với các phương pháp ước lượng ngưỡng như SURE hay BayesShrink, đã trở thành baseline mạnh và được nghiên cứu kỹ lưỡng cho bài toán này.

Tuy nhiên, DWT có một hạn chế cố hữu: nó chỉ phân tích ảnh theo ba hướng cố định --- ngang (horizontal), dọc (vertical) và chéo (diagonal). Đối với các loại ảnh có cấu trúc hướng phong phú như ảnh SAR (Synthetic Aperture Radar), ảnh y tế, hay các ảnh chứa nhiều texture có hướng xác định, biểu diễn bằng DWT không đủ sparse để khai thác triệt để sự khác biệt giữa signal và noise.

Để khắc phục giới hạn này, Bamberger và Smith \cite{bamberger1992} đề xuất Directional Filter Bank (DFB) --- một loại filter bank 2D có khả năng phân rã ảnh thành $N$ subband theo $N$ hướng tùy chọn trong mặt phẳng tần số, tạo thành các vùng hình nêm (wedge-shaped regions). Mỗi directional subband tập trung năng lượng của các cạnh và texture theo đúng hướng đó, giúp tín hiệu ảnh có biểu diễn sparse hơn trong những subband phù hợp với cấu trúc hướng của ảnh.

Dựa trên nền tảng này, Rosiles và Smith \cite{rosiles2000} đề xuất sử dụng DFB cho bài toán image denoising: áp dụng DFB 8 band để phân rã ảnh thành 8 directional subband, ước lượng threshold cho từng subband theo các tiêu chí thống kê như SURE (Stein's Unbiased Risk Estimate) hoặc GCV (Generalised Cross Validation), sau đó áp dụng soft thresholding và tái tạo ảnh bằng inverse DFB. Kết quả cho thấy DFB outperform DWT, đặc biệt trên các ảnh SAR có nhiều directional content.

Bài báo cáo này trình bày quá trình reimplementation thuật toán của Rosiles và Smith \cite{rosiles2000}, kèm theo một số cải tiến được đề xuất. Do tính phức tạp của việc cài đặt critically-sampled DFB với quincunx decimation như trong paper gốc, chúng tôi sử dụng một undecimated DFB trong miền tần số, kết hợp với kỹ thuật tách LP/HP để đảm bảo các directional subband có tính sparse phù hợp cho thresholding. Ngoài ra, ba cải tiến được đề xuất bao gồm: sử dụng BayesShrink thay cho SURE/GCV, tăng số band lên 16 để có độ phân giải góc mịn hơn, và phát triển một multi-scale undecimated DFB kết hợp Gaussian pyramid với DFB tại mỗi scale --- một kiến trúc có liên hệ chặt chẽ với Contourlet transform \cite{do2005}.

Phần còn lại của báo cáo được tổ chức như sau: Phần~\ref{sec:background} trình bày các kiến thức nền tảng về DFB và các phương pháp ước lượng threshold. Phần~\ref{sec:reimplementation} mô tả chi tiết quá trình reimplementation. Phần~\ref{sec:improvements} trình bày các cải tiến đề xuất. Phần~\ref{sec:experiments} và~\ref{sec:discussion} lần lượt trình bày kết quả thực nghiệm và thảo luận. Phần~\ref{sec:conclusion} kết luận báo cáo.
